%% kapa: exam=1 round=23 year=2012 period=2 form=A subject=appraisal_law subject_ko=부동산관계법규 seq=21 number=61
%% source: Q-Net 감정평가사 기출문제 (cst00309) artlSeq=5052580 file=raw/1008/5052580/2012년 제23회 감정평가사 제1차시험 기출문제(2교시 A형).pdf
%% conversion: kordoc 4.24.51 -> tools/gen1.py (auto); review=proofread-sonnet; flags=none
\begin{question}{61}
부동산 가격공시 및 감정평가에 관한 법령상 주택가격의 공시에 관한 설명으로 옳은 것은?
\begin{choices}
  \item 표준주택가격은 국토해양부장관이 중앙부동산평가위원회의 심의를 거쳐 조사ㆍ평가한 후 공시하여야 한다.
  \item 신축된 단독주택에 대하여 개별주택가격을 결정ㆍ공시할 경우에는 당해 연도의 공시기준일을 기준으로 한다.
  \item 165\,m$^2$ 미만인 연립주택의 공동주택가격은 국세청장이 결정ㆍ고시하여야 한다.
  \item 표준주택가격은 국가ㆍ지방자치단체 등의 기관이 그 업무와 관련하여 개별주택가격을 산정하는 경우에 그 기준이 된다.
  \item 개별주택 및 공동주택의 가격은 국가ㆍ지방자치단체 등의 기관이 수용방식의 도시계획사업을 시행할 때 보상가격으로 사용된다.
\end{choices}
\end{question}
\answer{4}
